\section{From one protein sequence to an alignment}
\label{sec:book-sequence-msa-primer}

A protein sequence records the amino-acid residues of a polymer chain in their order from one
terminus to the other.  A \emph{residue} is one amino-acid unit in that chain; a \emph{site} or
\emph{position} is the coordinate at which its identity is recorded.  The \emph{query} is the
particular sequence whose structure, function, or model response is under study.  A
\emph{homologue} is a sequence related to the query by common ancestry.  Sequence similarity is
evidence used to find homologues; homology itself is a historical relation.

A multiple-sequence alignment (MSA) places homologous sequences in rows and proposes
corresponding ancestral positions in columns.  The following toy alignment uses the query as its
coordinate system:
\begin{center}
\small
\setlength{\tabcolsep}{5pt}
\begin{tabular}{@{}lcccccccc@{}}
\toprule
row & 1 & 2 & 3 & 4 & 5 & 6 & 7 & 8 \\
\midrule
query       & M & K & T & A & E & L & V & G \\
homologue 1 & M & K & S & A & E & L & I & G \\
homologue 2 & M & R & T & A & D & L & V & G \\
homologue 3 & M & R & S & -- & D & L & I & G \\
\bottomrule
\end{tabular}
\end{center}
Columns 1, 6, and 8 are conserved in this sample.  Columns 2 and 5 covary: the observed rows pair
K with E and R with D.  The dash in column 4 is an alignment gap, representing an insertion,
deletion, missing observation, or an alignment decision according to the data pipeline.  It is not
an amino acid.  Row order in an alignment file has no biological meaning; relationships among
rows are supplied by common ancestry, taxonomy, sampling, or an inferred tree.

The model alphabet must therefore be declared.  The twenty standard amino acids may be augmented
with gap, unknown, ambiguous, or masked symbols.  In a query-anchored analysis, indices $i$ and
$j$ usually denote retained query-aligned columns.  They need not equal the raw residue numbers in
every homologue, and poorly aligned or insertion-rich regions can make the correspondence
uncertain.

Conservation is a one-column statement: the marginal distribution $p_i(a)$ is concentrated on a
small set of residue identities.  Covariation is a two-column statement: the joint distribution
$p_{ij}(a,b)$ differs from $p_i(a)p_j(b)$.  The toy alignment displays both, yet its four rows are
far too few to establish either quantity reliably.  More generally, observed covariation can
contain several mechanisms at once:
\begin{center}
\small
\begin{tabular}{@{}>{\raggedright\arraybackslash}p{0.24\linewidth}
                >{\raggedright\arraybackslash}p{0.36\linewidth}
                >{\raggedright\arraybackslash}p{0.30\linewidth}@{}}
\toprule
Source & How the pattern arises & What must be checked \\
\midrule
Direct structural or functional constraint & substitutions at two sites compensate to preserve
folding, binding, catalysis, or another selected property & matched structure, function, or
mutational evidence \\
Indirect dependence & $i$ and $j$ are both coupled to a third site $k$, or are joined by a longer
path & conditioning, inverse modeling, and controlled synthetic cases \\
Shared ancestry & related rows inherit both residue states from common ancestors & a tree-based or
clade-aware sampling null \\
Mixture and data construction & subfamilies, taxonomic imbalance, gaps, alignment choices, or
paired-MSA assignment change the empirical joint table & stratification and pipeline-matched nulls
\\
\bottomrule
\end{tabular}
\end{center}

Finite depth and database redundancy add uncertainty to every row of this ledger.  Sequence
weighting can reduce the influence of nearly duplicated observations, while the ancestral process
and alignment uncertainty remain part of the data-generating law.  An MSA is therefore an observed
three-axis object---row, aligned position, and residue category---from which a position-level
interaction operator may be estimated.  The estimator and its null determine which part of the
observed covariation the resulting edge can claim to represent.

With these three axes fixed, the next step is to distinguish operations performed on the resulting
many-variable system.  A global energy law, a factor-graph cavity message, a dynamical response,
and a cascade on a support graph act on different state spaces even when each is described with the
word ``interaction.''  The four statistical-mechanics lenses that follow provide exact reference
cases for those operations before the argument crosses from sequence data to physical and
evolutionary scales.
